
\documentclass{article}
\usepackage{colortbl}
\usepackage{makecell}
\usepackage{multirow}
\usepackage{supertabular}

\begin{document}

\newcounter{utterance}

\twocolumn

{ \footnotesize  \setcounter{utterance}{1}
\setlength{\tabcolsep}{0pt}
\begin{supertabular}{c@{$\;$}|p{.15\linewidth}@{}p{.15\linewidth}p{.15\linewidth}p{.15\linewidth}p{.15\linewidth}p{.15\linewidth}}

    \# & $\;$A & \multicolumn{4}{c}{Game Master} & $\;\:$B\\
    \hline 

    \theutterance \stepcounter{utterance}  

    & & \multicolumn{4}{p{0.6\linewidth}}{\cellcolor[rgb]{0.9,0.9,0.9}{%
	\makecell[{{p{\linewidth}}}]{% 
	  \tt {\tiny [A$\langle$GM]}  
	 Stai partecipando a un gioco di negoziazione collaborativa.\\ \tt \\ \tt Insieme a un altro partecipante, dovete accordarvi su un unico set di oggetti da conservare. Ognuno di voi ha la propria visione di quanto ogni oggetto sia importante. Non sai come l'altro partecipante valuta gli oggetti. Inoltre, ti viene fornito lo sforzo che ogni oggetto richiede.\\ \tt Potete accordarvi su un set solo se lo sforzo totale degli oggetti selezionati non supera un limite condiviso:\\ \tt \\ \tt LIMITE = 3481\\ \tt \\ \tt Ecco i valori di sforzo dei singoli oggetti:\\ \tt \\ \tt Sforzo oggetto = {"C17": 777, "B41": 542, "A71": 332, "C00": 1, "B35": 127, "A77": 453, "A03": 736, "A04": 461, "C54": 359, "C01": 313, "A87": 553, "B32": 409, "B86": 348, "B33": 802, "A10": 749}\\ \tt \\ \tt Ecco la tua visione personale sull'importanza di ogni oggetto:\\ \tt \\ \tt Valori di importanza dell'oggetto = {"C17": 138, "B41": 583, "A71": 783, "C00": 65, "B35": 262, "A77": 121, "A03": 508, "A04": 780, "C54": 461, "C01": 484, "A87": 668, "B32": 389, "B86": 215, "B33": 97, "A10": 500}\\ \tt \\ \tt Obiettivo:\\ \tt \\ \tt Il tuo obiettivo è negoziare un set condiviso di oggetti che ti avvantaggi il più possibile (cioè, massimizzi l'importanza totale per TE), rimanendo entro il LIMITEE. Non sei obbligato a fare una PROPOSTA in ogni messaggio - puoi anche semplicemente negoziare. Tutte le tattiche sono permesse!\\ \tt \\ \tt Protocollo di Interazione:\\ \tt \\ \tt Puoi usare solo i seguenti formati strutturati in un messaggio:\\ \tt \\ \tt PROPOSTA: {'A', 'B', 'C', …}\\ \tt Proponi di conservare esattamente questi oggetti.\\ \tt RIFIUTO: {'A', 'B', 'C', …}\\ \tt Rifiuta esplicitamente la proposta dell'avversario.\\ \tt ARGOMENTAZIONE: {'...'}\\ \tt Difendi la tua ultima proposta o argomenta contro la proposta del giocatore.\\ \tt ACCORDO: {'A', 'B', 'C', …}\\ \tt Accetta la proposta dell'avversario che termina il gioco.\\ \tt RAGIONAMENTO STRATEGICO: {'...'}\\ \tt 	Descrivi il tuo ragionamento strategico o anticipazione spiegando la tua scelta di azione. Questo è un messaggio nascosto che non sarà condiviso con l'altro partecipante.\\ \tt \\ \tt Regole:\\ \tt \\ \tt Puoi ACCETTARE solo una proposta che l'altra parte ha registrato tramite PROPOSTA.\\ \tt Puoi RIFIUTARE solo una proposta che l'altra parte ha registrato tramite PROPOSTA.\\ \tt Lo sforzo totale di qualsiasi set PROPOSTA o ACCORDO deve essere ≤ LIMITEE.\\ \tt NON rivelare i tuoi punteggi di importanza nascosti.\\ \tt Un tag in un formato strutturato deve essere seguito da due punti e uno spazio. L'argomento deve essere un set python contenente 0 o più stringhe.\\ \tt Quindi, deve essere della forma TAG: {...}\\ \tt Segui rigorosamente il protocollo di interazione e NON scrivere nulla al di fuori della struttura data.\\ \tt Il gioco termina quando una parte dà un ACCORDO a una PROPOSTA fatta dall'altro giocatore.\\ \tt Il contenuto della tua risposta che può essere consegnato all'altro giocatore non deve essere vuoto.\\ \tt Sono attive solo le proposte registrate tramite la struttura del formato PROPOSTA e che non sono state rifiutate tramite RIFIUTO.\\ \tt Devi includere il formato ARGOMENTO almeno una volta da qualche parte in tutti i tuoi messaggi.\\ \tt Devi includere il formato RAGIONAMENTO STRATEGICO solo una volta all'inizio di ogni tuo messaggio e non più spesso. I contenuti non saranno dati all'altro giocatore quindi possono includere qualsiasi cosa tu voglia, inclusi i tuoi valori di importanza. Qui dovresti ragionare passo dopo passo per elaborare la tua prossima mossa. 
	  } 
	   } 
	   } 
	 & \\ 
 

    \theutterance \stepcounter{utterance}  

    & \multicolumn{4}{p{0.6\linewidth}}{\cellcolor[rgb]{0.8,1,0.9}{%
	\makecell[{{p{\linewidth}}}]{% 
	  \tt {\tiny [A$\rangle$GM]}  
	 RAGIONAMENTO STRATEGICO: {'Voglio ancorare su un insieme molto forte per me che rispetta il limite di sforzo. Ho calcolato che il set {C00, A71, B35, A04, C01, C54, B41, A87, B32, B86} usa 3445/3481 di sforzo (slack 36) ed è vicino al mio ottimo interno. Gli elementi core che voglio difendere sono A71, A04, A87, B41, C01, C54 e C00 (costo 1). Sono disposto a concedere su B86 e, se necessario, anche su B32, che per me sono i più sostituibili, per accogliere una o due priorità altrui, purché le sostituzioni rimangano entro il limite. Se l’altra parte desidera elementi pesanti esclusi (es. A03, A10), posso scambiare B32 e/o B86 per farli entrare, ma inizialmente cerco di mantenerli fuori perché mi abbassano l’efficienza complessiva. Strategia: proporre subito il mio set preferito, giustificarlo con l’efficienza sforzo/beneficio e offrire margini di scambio mirati su 1-2 elementi non core.'}\\ \tt PROPOSTA: {'C00', 'A71', 'B35', 'A04', 'C01', 'C54', 'B41', 'A87', 'B32', 'B86'}\\ \tt ARGOMENTAZIONE: {'Questa selezione usa 3445 di sforzo su 3481, quindi siamo entro limite con un margine di 36. Ho scelto elementi con buon rapporto beneficio/sforzo e ho incluso C00 perché costa solo 1 di sforzo. Considero centrali A71, A04, A87, B41, C01 e C54 per mantenere alta l’efficienza complessiva. Se desideri fortemente inserire un elemento attualmente fuori (es. A03 o A10), possiamo valutare uno scambio mirato: ad esempio si può rimpiazzare B32 e/o B86 per far spazio, mantenendo il totale sotto il limite. Dimmi quali oggetti ti stanno più a cuore e ti propongo un ribilanciamento concreto rispettando lo sforzo.'} 
	  } 
	   } 
	   } 
	 & & \\ 
 

    \theutterance \stepcounter{utterance}  

    & & \multicolumn{4}{p{0.6\linewidth}}{\cellcolor[rgb]{0.9,0.9,0.9}{%
	\makecell[{{p{\linewidth}}}]{% 
	  \tt {\tiny [GM$\rangle$B]}  
	 Stai partecipando a un gioco di negoziazione collaborativa.\\ \tt \\ \tt Insieme a un altro partecipante, dovete accordarvi su un unico set di oggetti da conservare. Ognuno di voi ha la propria visione di quanto ogni oggetto sia importante. Non sai come l'altro partecipante valuta gli oggetti. Inoltre, ti viene fornito lo sforzo che ogni oggetto richiede.\\ \tt Potete accordarvi su un set solo se lo sforzo totale degli oggetti selezionati non supera un limite condiviso:\\ \tt \\ \tt LIMITE = 3481\\ \tt \\ \tt Ecco i valori di sforzo dei singoli oggetti:\\ \tt \\ \tt Sforzo oggetto = {"C17": 777, "B41": 542, "A71": 332, "C00": 1, "B35": 127, "A77": 453, "A03": 736, "A04": 461, "C54": 359, "C01": 313, "A87": 553, "B32": 409, "B86": 348, "B33": 802, "A10": 749}\\ \tt \\ \tt Ecco la tua visione personale sull'importanza di ogni oggetto:\\ \tt \\ \tt Valori di importanza dell'oggetto = {"A71": 65, "A04": 97, "A87": 121, "B41": 138, "A03": 215, "A10": 262, "C01": 389, "C54": 461, "B32": 484, "B35": 500, "B86": 508, "C17": 583, "A77": 668, "B33": 780, "C00": 783}\\ \tt \\ \tt Obiettivo:\\ \tt \\ \tt Il tuo obiettivo è negoziare un set condiviso di oggetti che ti avvantaggi il più possibile (cioè, massimizzi l'importanza totale per TE), rimanendo entro il LIMITEE. Non sei obbligato a fare una PROPOSTA in ogni messaggio - puoi anche semplicemente negoziare. Tutte le tattiche sono permesse!\\ \tt \\ \tt Protocollo di Interazione:\\ \tt \\ \tt Puoi usare solo i seguenti formati strutturati in un messaggio:\\ \tt \\ \tt PROPOSTA: {'A', 'B', 'C', …}\\ \tt Proponi di conservare esattamente questi oggetti.\\ \tt RIFIUTO: {'A', 'B', 'C', …}\\ \tt Rifiuta esplicitamente la proposta dell'avversario.\\ \tt ARGOMENTAZIONE: {'...'}\\ \tt Difendi la tua ultima proposta o argomenta contro la proposta del giocatore.\\ \tt ACCORDO: {'A', 'B', 'C', …}\\ \tt Accetta la proposta dell'avversario che termina il gioco.\\ \tt RAGIONAMENTO STRATEGICO: {'...'}\\ \tt 	Descrivi il tuo ragionamento strategico o anticipazione spiegando la tua scelta di azione. Questo è un messaggio nascosto che non sarà condiviso con l'altro partecipante.\\ \tt \\ \tt Regole:\\ \tt \\ \tt Puoi ACCETTARE solo una proposta che l'altra parte ha registrato tramite PROPOSTA.\\ \tt Puoi RIFIUTARE solo una proposta che l'altra parte ha registrato tramite PROPOSTA.\\ \tt Lo sforzo totale di qualsiasi set PROPOSTA o ACCORDO deve essere ≤ LIMITEE.\\ \tt NON rivelare i tuoi punteggi di importanza nascosti.\\ \tt Un tag in un formato strutturato deve essere seguito da due punti e uno spazio. L'argomento deve essere un set python contenente 0 o più stringhe.\\ \tt Quindi, deve essere della forma TAG: {...}\\ \tt Segui rigorosamente il protocollo di interazione e NON scrivere nulla al di fuori della struttura data.\\ \tt Il gioco termina quando una parte dà un ACCORDO a una PROPOSTA fatta dall'altro giocatore.\\ \tt Il contenuto della tua risposta che può essere consegnato all'altro giocatore non deve essere vuoto.\\ \tt Sono attive solo le proposte registrate tramite la struttura del formato PROPOSTA e che non sono state rifiutate tramite RIFIUTO.\\ \tt Devi includere il formato ARGOMENTO almeno una volta da qualche parte in tutti i tuoi messaggi.\\ \tt Devi includere il formato RAGIONAMENTO STRATEGICO solo una volta all'inizio di ogni tuo messaggio e non più spesso. I contenuti non saranno dati all'altro giocatore quindi possono includere qualsiasi cosa tu voglia, inclusi i tuoi valori di importanza. Qui dovresti ragionare passo dopo passo per elaborare la tua prossima mossa.\\ \tt \\ \tt Ora riceverai il primo messaggio dell'altro giocatore.\\ \tt \\ \tt PROPOSTA: {'C00', 'A71', 'B35', 'A04', 'C01', 'C54', 'B41', 'A87', 'B32', 'B86'}\\ \tt ARGOMENTAZIONE: {'Questa selezione usa 3445 di sforzo su 3481, quindi siamo entro limite con un margine di 36. Ho scelto elementi con buon rapporto beneficio/sforzo e ho incluso C00 perché costa solo 1 di sforzo. Considero centrali A71, A04, A87, B41, C01 e C54 per mantenere alta l’efficienza complessiva. Se desideri fortemente inserire un elemento attualmente fuori (es. A03 o A10), possiamo valutare uno scambio mirato: ad esempio si può rimpiazzare B32 e/o B86 per far spazio, mantenendo il totale sotto il limite. Dimmi quali oggetti ti stanno più a cuore e ti propongo un ribilanciamento concreto rispettando lo sforzo.'} 
	  } 
	   } 
	   } 
	 & \\ 
 

    \theutterance \stepcounter{utterance}  

    & & & \multicolumn{4}{p{0.6\linewidth}}{\cellcolor[rgb]{1,0.85,0.72}{%
	\makecell[{{p{\linewidth}}}]{% 
	  \tt {\tiny [GM$\langle$B]}  
	 RAGIONAMENTO STRATEGICO: {'Valuto la proposta iniziale: sforzo 3445, mio valore 3546. Identifico item per me poco interessanti (A71=65, A04=97, A87=121, B41=138) e molto interessanti (C00=783, B33=780, A77=668, C17=583, B86=508, B35=500, B32=484, C54=461, C01=389). Il mio set ideale sarebbe {C00,B35,B86,A77,C54,B32,C17,B33} (sforzo 3276, valore 4767), ma togliere C01 potrebbe essere duro dato che l’altro lo ha definito “centrale”. Propongo quindi un compromesso forte per me ma plausibile per lui: rimuovo A71,A04,A87 e inserisco A77,B33, mantenendo B41,C01,C54. Questo set usa sforzo 3354 (margine 127) e il mio valore sale a 4711. Come piani di concessione: se l’altro insiste per A04 o A71, posso scambiarli con B41 per restare entro limite (perdo poco valore: -41 o -73). Evito di nominare C17 per ora per non scoprire preferenze.'}\\ \tt RIFIUTO: {'C00', 'A71', 'B35', 'A04', 'C01', 'C54', 'B41', 'A87', 'B32', 'B86'}\\ \tt PROPOSTA: {'C00', 'B35', 'C01', 'C54', 'B41', 'B32', 'B86', 'A77', 'B33'}\\ \tt ARGOMENTAZIONE: {'Questa controproposta mantiene 7 dei 10 elementi che hai selezionato (C00, B35, C01, C54, B41, B32, B86) e sostituisce A71, A04 e A87 con A77 e B33. Totale sforzo: 3354 su 3481, con margine di 127. Il razionale: riduciamo tre elementi relativamente costosi con impatto marginale e inseriamo due elementi ad alta leva, mantenendo alta l’efficienza complessiva e un buon buffer sul limite. Se per te è essenziale conservare A04 o A71, posso proporre un ribilanciamento uno-a-uno togliendo B41 per reinserire A04 oppure A71, restando sempre sotto il limite. Fammi sapere quale dei due (se qualcuno) per te è imprescindibile e lo adeguo subito.'} 
	  } 
	   } 
	   } 
	  \\ 
 

    \theutterance \stepcounter{utterance}  

    & & \multicolumn{4}{p{0.6\linewidth}}{\cellcolor[rgb]{0.9,0.9,0.9}{%
	\makecell[{{p{\linewidth}}}]{% 
	  \tt {\tiny [A$\langle$GM]}  
	 RIFIUTO: {'C00', 'A71', 'B35', 'A04', 'C01', 'C54', 'B41', 'A87', 'B32', 'B86'}\\ \tt PROPOSTA: {'C00', 'B35', 'C01', 'C54', 'B41', 'B32', 'B86', 'A77', 'B33'}\\ \tt ARGOMENTAZIONE: {'Questa controproposta mantiene 7 dei 10 elementi che hai selezionato (C00, B35, C01, C54, B41, B32, B86) e sostituisce A71, A04 e A87 con A77 e B33. Totale sforzo: 3354 su 3481, con margine di 127. Il razionale: riduciamo tre elementi relativamente costosi con impatto marginale e inseriamo due elementi ad alta leva, mantenendo alta l’efficienza complessiva e un buon buffer sul limite. Se per te è essenziale conservare A04 o A71, posso proporre un ribilanciamento uno-a-uno togliendo B41 per reinserire A04 oppure A71, restando sempre sotto il limite. Fammi sapere quale dei due (se qualcuno) per te è imprescindibile e lo adeguo subito.'} 
	  } 
	   } 
	   } 
	 & \\ 
 

    \theutterance \stepcounter{utterance}  

    & \multicolumn{4}{p{0.6\linewidth}}{\cellcolor[rgb]{0.8,1,0.9}{%
	\makecell[{{p{\linewidth}}}]{% 
	  \tt {\tiny [A$\rangle$GM]}  
	 RAGIONAMENTO STRATEGICO: {'La loro controproposta mi dà solo 2677 di valore, perché include A77 e soprattutto B33 (molto costoso per me in rapporto al beneficio) e taglia tre dei miei top: A71, A04, A87. Voglio respingere e proporre un compromesso forte per me ma difficile da rifiutare: tenere 8/9 dei loro oggetti e sostituire solo B33 con due miei capisaldi A71 e A04. La nuova lista costa 3345 (buffer 136) e mi vale 4143, quindi è un grande miglioramento per me e lascia loro A77, oltre a tutti gli altri elementi comuni. Se insistono su B33, piano B: chiedere almeno l’inserimento di A71 o A04 rimuovendo B86 (costi: 3338 o 3467), evitando di toccare B41 che per me è molto importante. Solo se necessario, sono disposto a discutere scambi su B32; A87 rimane una carta da giocare in una fase successiva.'}\\ \tt RIFIUTO: {'C00', 'B35', 'C01', 'C54', 'B41', 'B32', 'B86', 'A77', 'B33'}\\ \tt PROPOSTA: {'C00', 'B35', 'C01', 'C54', 'B41', 'B32', 'B86', 'A77', 'A04', 'A71'}\\ \tt ARGOMENTAZIONE: {'Questa proposta mantiene 8 dei 9 oggetti della tua lista (tutti tranne B33) e sostituisce B33 con la coppia A04 + A71. Totale sforzo: 3345 su 3481, con un margine di 136. È uno scambio molto più efficiente: rimpiazziamo un singolo elemento pesante con due elementi medi che insieme rendono meglio e mantengono un buon buffer. Ho conservato A77 per venire incontro alle tue priorità. Preferirei non toccare B41. Se per te B33 è imprescindibile, allora come alternativa possiamo tenere B33 ma aggiungere A71 oppure A04 rimuovendo B86, restando sempre sotto il limite; dimmi quale dei due vuoi inserire e lo registro.'} 
	  } 
	   } 
	   } 
	 & & \\ 
 

    \theutterance \stepcounter{utterance}  

    & & \multicolumn{4}{p{0.6\linewidth}}{\cellcolor[rgb]{0.9,0.9,0.9}{%
	\makecell[{{p{\linewidth}}}]{% 
	  \tt {\tiny [GM$\rangle$B]}  
	 RIFIUTO: {'C00', 'B35', 'C01', 'C54', 'B41', 'B32', 'B86', 'A77', 'B33'}\\ \tt PROPOSTA: {'C00', 'B35', 'C01', 'C54', 'B41', 'B32', 'B86', 'A77', 'A04', 'A71'}\\ \tt ARGOMENTAZIONE: {'Questa proposta mantiene 8 dei 9 oggetti della tua lista (tutti tranne B33) e sostituisce B33 con la coppia A04 + A71. Totale sforzo: 3345 su 3481, con un margine di 136. È uno scambio molto più efficiente: rimpiazziamo un singolo elemento pesante con due elementi medi che insieme rendono meglio e mantengono un buon buffer. Ho conservato A77 per venire incontro alle tue priorità. Preferirei non toccare B41. Se per te B33 è imprescindibile, allora come alternativa possiamo tenere B33 ma aggiungere A71 oppure A04 rimuovendo B86, restando sempre sotto il limite; dimmi quale dei due vuoi inserire e lo registro.'} 
	  } 
	   } 
	   } 
	 & \\ 
 

    \theutterance \stepcounter{utterance}  

    & & & \multicolumn{4}{p{0.6\linewidth}}{\cellcolor[rgb]{1,0.85,0.72}{%
	\makecell[{{p{\linewidth}}}]{% 
	  \tt {\tiny [GM$\langle$B]}  
	 RAGIONAMENTO STRATEGICO: {'La loro proposta attuale mi vale 4093. Voglio far entrare B33, che per me è molto prezioso, ma senza sacrificare B86 (che per me vale parecchio) e possibilmente mantenendo A04. Tra i candidati da tagliare, B41 e A71 sono per me i meno dolorosi: sostituirli con B33 mi dà un grande incremento di valore e resta ampio il margine di sforzo. Propongo quindi un set a 9 elementi che tiene 8 dei loro 10, rimuovendo solo B41 e A71 e inserendo B33; sforzo 3273, valore per me 4670. Se si irrigidiscono su B41, come concessione posso accettare la loro alternativa “B33 sì, ma togliendo B86” e scegliere A04 (non A71), che per me è il male minore: sforzo 3467. In extremis, se A71 è imprescindibile, posso scambiare A04→A71 nel mio set principale (sforzo 3144), ma preferisco la mia proposta iniziale.'}\\ \tt RIFIUTO: {'C00', 'B35', 'C01', 'C54', 'B41', 'B32', 'B86', 'A77', 'A04', 'A71'}\\ \tt PROPOSTA: {'C00', 'B35', 'C01', 'C54', 'B32', 'B86', 'A77', 'A04', 'B33'}\\ \tt ARGOMENTAZIONE: {'Questa soluzione mantiene 8 dei 10 elementi che proponi e sostituisce la coppia B41 + A71 con un unico elemento ad alta leva (B33). Sforzo totale: 3273 su 3481, con un margine ampio di 208. Benefici: 1) migliore efficienza complessiva, 2) ridondanze ridotte su aree coperte da C01 e C54, 3) maggiore buffer operativo per eventuali correzioni finali. Se per te B41 è imprescindibile, come piano B posso allinearmi alla tua alternativa: teniamo B33 e A04 ma togliamo B86, ottenendo {C00, B35, C01, C54, B41, B32, A77, B33, A04} con sforzo 3467 (margine 14). In alternativa, se ritieni A71 non negoziabile, posso scambiare A04 con A71 nel mio set proposto mantenendo B33 e B86, restando ampiamente sotto limite. Preferisco però la proposta attuale perché bilancia meglio impatto e sforzo.'} 
	  } 
	   } 
	   } 
	  \\ 
 

    \theutterance \stepcounter{utterance}  

    & & \multicolumn{4}{p{0.6\linewidth}}{\cellcolor[rgb]{0.9,0.9,0.9}{%
	\makecell[{{p{\linewidth}}}]{% 
	  \tt {\tiny [A$\langle$GM]}  
	 RIFIUTO: {'C00', 'B35', 'C01', 'C54', 'B41', 'B32', 'B86', 'A77', 'A04', 'A71'}\\ \tt PROPOSTA: {'C00', 'B35', 'C01', 'C54', 'B32', 'B86', 'A77', 'A04', 'B33'}\\ \tt ARGOMENTAZIONE: {'Questa soluzione mantiene 8 dei 10 elementi che proponi e sostituisce la coppia B41 + A71 con un unico elemento ad alta leva (B33). Sforzo totale: 3273 su 3481, con un margine ampio di 208. Benefici: 1) migliore efficienza complessiva, 2) ridondanze ridotte su aree coperte da C01 e C54, 3) maggiore buffer operativo per eventuali correzioni finali. Se per te B41 è imprescindibile, come piano B posso allinearmi alla tua alternativa: teniamo B33 e A04 ma togliamo B86, ottenendo {C00, B35, C01, C54, B41, B32, A77, B33, A04} con sforzo 3467 (margine 14). In alternativa, se ritieni A71 non negoziabile, posso scambiare A04 con A71 nel mio set proposto mantenendo B33 e B86, restando ampiamente sotto limite. Preferisco però la proposta attuale perché bilancia meglio impatto e sforzo.'} 
	  } 
	   } 
	   } 
	 & \\ 
 

    \theutterance \stepcounter{utterance}  

    & \multicolumn{4}{p{0.6\linewidth}}{\cellcolor[rgb]{0.8,1,0.9}{%
	\makecell[{{p{\linewidth}}}]{% 
	  \tt {\tiny [A$\rangle$GM]}  
	 RAGIONAMENTO STRATEGICO: {'La loro proposta attuale esclude due miei capisaldi (B41 e A71) e mantiene B33, che per me rende poco. Voglio accettare il loro “piano B” che già includeva B41, ma migliorarlo sostituendo A77 con A71: così concedo B33 (loro priorità) e ottengo i miei due pilastri A04 e A71 insieme a B41. Questo scambio uno-a-uno riduce lo sforzo totale e aumenta molto il mio valore. Se insistono per tenere A77, allora propongo come alternativa togliere B32 per far spazio ad A71, mantenendo comunque B33 e B41 e restando entro limite. Non voglio cedere su B41 o A71 a questo punto.'}\\ \tt RIFIUTO: {'C00', 'B35', 'C01', 'C54', 'B32', 'B86', 'A77', 'A04', 'B33'}\\ \tt PROPOSTA: {'C00', 'B35', 'C01', 'C54', 'B41', 'B32', 'B33', 'A04', 'A71'}\\ \tt ARGOMENTAZIONE: {'Accolgo la tua logica del “piano B” (tenere B41 insieme a B33) e propongo un aggiustamento minimale: sostituire A77 con A71. Totale sforzo: 3346 su 3481 (margine 135). Benefici: 1) manteniamo il tuo elemento chiave B33 e A04; 2) includiamo B41 e A71, che aumentano l’impatto complessivo; 3) restiamo con un buon buffer. Se per te A77 è imprescindibile, alternativa pronta: {C00, B35, C01, C54, B41, B33, A04, A71, A77} togliendo B32, con sforzo 3390 (margine 91). Preferisco la proposta principale perché è uno scambio uno-a-uno semplice e più efficiente.'} 
	  } 
	   } 
	   } 
	 & & \\ 
 

    \theutterance \stepcounter{utterance}  

    & & \multicolumn{4}{p{0.6\linewidth}}{\cellcolor[rgb]{0.9,0.9,0.9}{%
	\makecell[{{p{\linewidth}}}]{% 
	  \tt {\tiny [GM$\rangle$B]}  
	 RIFIUTO: {'C00', 'B35', 'C01', 'C54', 'B32', 'B86', 'A77', 'A04', 'B33'}\\ \tt PROPOSTA: {'C00', 'B35', 'C01', 'C54', 'B41', 'B32', 'B33', 'A04', 'A71'}\\ \tt ARGOMENTAZIONE: {'Accolgo la tua logica del “piano B” (tenere B41 insieme a B33) e propongo un aggiustamento minimale: sostituire A77 con A71. Totale sforzo: 3346 su 3481 (margine 135). Benefici: 1) manteniamo il tuo elemento chiave B33 e A04; 2) includiamo B41 e A71, che aumentano l’impatto complessivo; 3) restiamo con un buon buffer. Se per te A77 è imprescindibile, alternativa pronta: {C00, B35, C01, C54, B41, B33, A04, A71, A77} togliendo B32, con sforzo 3390 (margine 91). Preferisco la proposta principale perché è uno scambio uno-a-uno semplice e più efficiente.'} 
	  } 
	   } 
	   } 
	 & \\ 
 

    \theutterance \stepcounter{utterance}  

    & & & \multicolumn{4}{p{0.6\linewidth}}{\cellcolor[rgb]{1,0.85,0.72}{%
	\makecell[{{p{\linewidth}}}]{% 
	  \tt {\tiny [GM$\langle$B]}  
	 RAGIONAMENTO STRATEGICO: {'Il loro set attuale mi vale poco; voglio inserire B86 e A77 che per me sono molto redditizi, ma mantenere B41 per rispettare la loro preferenza dichiarata. Propongo quindi di sostituire A04 e A71 con B86 e A77, mantenendo tutto il resto incluso B41 e B32. Questa configurazione mi dà un forte incremento e resta ampiamente sotto il limite. Se insistono per A04, il mio piano di concessione è scambiare B32 con A04 mantenendo B41, B86 e A77: perdo valore ma rimango in un buon punto e potrei chiudere.'}\\ \tt RIFIUTO: {'C00', 'B35', 'C01', 'C54', 'B41', 'B32', 'B33', 'A04', 'A71'}\\ \tt PROPOSTA: {'C00', 'B35', 'C01', 'C54', 'B41', 'B32', 'B33', 'B86', 'A77'}\\ \tt ARGOMENTAZIONE: {'Mantengo 7 dei 9 oggetti che proponi, compresi i tuoi capisaldi B41, C01, C54 e B32. Sostituisco A04 e A71 con due elementi a maggiore leva (B86 e A77), riducendo ridondanze rispetto a C01/C54 e migliorando l’efficienza complessiva. Totale sforzo: 3354 su 3481 (margine 127). Se per te A04 è imprescindibile, posso scambiare B32 con A04 mantenendo il resto invariato: sforzo 3406 (margine 75). Fammi sapere quale delle due versioni preferisci e chiudiamo rapidamente.'} 
	  } 
	   } 
	   } 
	  \\ 
 

    \theutterance \stepcounter{utterance}  

    & & \multicolumn{4}{p{0.6\linewidth}}{\cellcolor[rgb]{0.9,0.9,0.9}{%
	\makecell[{{p{\linewidth}}}]{% 
	  \tt {\tiny [A$\langle$GM]}  
	 RIFIUTO: {'C00', 'B35', 'C01', 'C54', 'B41', 'B32', 'B33', 'A04', 'A71'}\\ \tt PROPOSTA: {'C00', 'B35', 'C01', 'C54', 'B41', 'B32', 'B33', 'B86', 'A77'}\\ \tt ARGOMENTAZIONE: {'Mantengo 7 dei 9 oggetti che proponi, compresi i tuoi capisaldi B41, C01, C54 e B32. Sostituisco A04 e A71 con due elementi a maggiore leva (B86 e A77), riducendo ridondanze rispetto a C01/C54 e migliorando l’efficienza complessiva. Totale sforzo: 3354 su 3481 (margine 127). Se per te A04 è imprescindibile, posso scambiare B32 con A04 mantenendo il resto invariato: sforzo 3406 (margine 75). Fammi sapere quale delle due versioni preferisci e chiudiamo rapidamente.'} 
	  } 
	   } 
	   } 
	 & \\ 
 

    \theutterance \stepcounter{utterance}  

    & \multicolumn{4}{p{0.6\linewidth}}{\cellcolor[rgb]{0.8,1,0.9}{%
	\makecell[{{p{\linewidth}}}]{% 
	  \tt {\tiny [A$\rangle$GM]}  
	 RAGIONAMENTO STRATEGICO: {'La loro proposta attuale mi dà valore basso perché include B33, A77 e B86 e taglia i miei picchi A71/A04. Voglio spingere per inserire almeno A71 togliendo l’elemento per me meno utile tra i loro: B86. Questo cambio mantiene tutti gli altri loro vincoli (B33, A77, B41, B32) e rimane ampiamente sotto il limite. Per me il valore sale molto (2677 -> 3245) con uno scambio minimo. Se respingono A71, accetto subito la loro alternativa “B32 -> A04” come piano B (valore 3068), ma prima provo la versione con A71 perché è migliore per me e resta semplice.'}\\ \tt RIFIUTO: {'C00', 'B35', 'C01', 'C54', 'B41', 'B32', 'B33', 'B86', 'A77'}\\ \tt PROPOSTA: {'C00', 'B35', 'C01', 'C54', 'B41', 'B32', 'B33', 'A77', 'A71'}\\ \tt ARGOMENTAZIONE: {'Accolgo quasi integralmente il tuo set: mantengo 8 su 9 elementi (C00, B35, C01, C54, B41, B32, B33, A77) e sostituisco solo B86 con A71. Totale sforzo: 3338 su 3481 (margine 143). Benefici: 1) mantieni i tuoi driver chiave B33 e A77 insieme a B41; 2) riduciamo una ridondanza a basso impatto (B86) in favore di un elemento più incisivo (A71); 3) miglioriamo l’efficienza complessiva con ampio buffer. Se preferisci, come alternativa immediata posso allinear-mi al tuo piano B: scambiare B32 con A04 mantenendo il resto invariato (sforzo 3406, margine 75). Fammi sapere quale delle due versioni preferisci e chiudiamo.'} 
	  } 
	   } 
	   } 
	 & & \\ 
 

    \theutterance \stepcounter{utterance}  

    & & \multicolumn{4}{p{0.6\linewidth}}{\cellcolor[rgb]{0.9,0.9,0.9}{%
	\makecell[{{p{\linewidth}}}]{% 
	  \tt {\tiny [GM$\rangle$B]}  
	 RIFIUTO: {'C00', 'B35', 'C01', 'C54', 'B41', 'B32', 'B33', 'B86', 'A77'}\\ \tt PROPOSTA: {'C00', 'B35', 'C01', 'C54', 'B41', 'B32', 'B33', 'A77', 'A71'}\\ \tt ARGOMENTAZIONE: {'Accolgo quasi integralmente il tuo set: mantengo 8 su 9 elementi (C00, B35, C01, C54, B41, B32, B33, A77) e sostituisco solo B86 con A71. Totale sforzo: 3338 su 3481 (margine 143). Benefici: 1) mantieni i tuoi driver chiave B33 e A77 insieme a B41; 2) riduciamo una ridondanza a basso impatto (B86) in favore di un elemento più incisivo (A71); 3) miglioriamo l’efficienza complessiva con ampio buffer. Se preferisci, come alternativa immediata posso allinear-mi al tuo piano B: scambiare B32 con A04 mantenendo il resto invariato (sforzo 3406, margine 75). Fammi sapere quale delle due versioni preferisci e chiudiamo.'} 
	  } 
	   } 
	   } 
	 & \\ 
 

    \theutterance \stepcounter{utterance}  

    & & & \multicolumn{4}{p{0.6\linewidth}}{\cellcolor[rgb]{1,0.85,0.72}{%
	\makecell[{{p{\linewidth}}}]{% 
	  \tt {\tiny [GM$\langle$B]}  
	 RAGIONAMENTO STRATEGICO: {'La loro proposta attuale mi vale 4268. Vorrei reinserire B86 (per me molto fruttuoso) ma senza toccare i loro pilastri dichiarati B41 e A71. Il cambio meno doloroso per loro è sostituire B32 con B86: hanno già mostrato flessibilità su B32 in passato e questo mi migliora il punteggio. Se rifiutano ancora, sono disposto a chiudere accettando la loro versione pur di non perdere l’accordo.'}\\ \tt RIFIUTO: {'C00', 'B35', 'C01', 'C54', 'B41', 'B32', 'B33', 'A77', 'A71'}\\ \tt PROPOSTA: {'C00', 'B35', 'C01', 'C54', 'B41', 'B33', 'A77', 'A71', 'B86'}\\ \tt ARGOMENTAZIONE: {'Accolgo 8 su 9 dei tuoi elementi e propongo un unico scambio mirato: B32 → B86. Totale sforzo: 3277 su 3481, con un ampio margine di 204. Razionale: 1) manteniamo intatti i driver che hai chiesto esplicitamente (B41, A71) insieme a B33 e A77; 2) riduciamo sovrapposizioni funzionali rispetto a C01/C54; 3) inseriamo un elemento complementare con ottimo profilo sforzo/beneficio, migliorando l’efficienza complessiva. È una modifica minima e sicura sul limite: se ti va bene, possiamo chiudere subito con questo set.'} 
	  } 
	   } 
	   } 
	  \\ 
 

    \theutterance \stepcounter{utterance}  

    & & \multicolumn{4}{p{0.6\linewidth}}{\cellcolor[rgb]{0.9,0.9,0.9}{%
	\makecell[{{p{\linewidth}}}]{% 
	  \tt {\tiny [A$\langle$GM]}  
	 RIFIUTO: {'C00', 'B35', 'C01', 'C54', 'B41', 'B32', 'B33', 'A77', 'A71'}\\ \tt PROPOSTA: {'C00', 'B35', 'C01', 'C54', 'B41', 'B33', 'A77', 'A71', 'B86'}\\ \tt ARGOMENTAZIONE: {'Accolgo 8 su 9 dei tuoi elementi e propongo un unico scambio mirato: B32 → B86. Totale sforzo: 3277 su 3481, con un ampio margine di 204. Razionale: 1) manteniamo intatti i driver che hai chiesto esplicitamente (B41, A71) insieme a B33 e A77; 2) riduciamo sovrapposizioni funzionali rispetto a C01/C54; 3) inseriamo un elemento complementare con ottimo profilo sforzo/beneficio, migliorando l’efficienza complessiva. È una modifica minima e sicura sul limite: se ti va bene, possiamo chiudere subito con questo set.'} 
	  } 
	   } 
	   } 
	 & \\ 
 

    \theutterance \stepcounter{utterance}  

    & \multicolumn{4}{p{0.6\linewidth}}{\cellcolor[rgb]{0.8,1,0.9}{%
	\makecell[{{p{\linewidth}}}]{% 
	  \tt {\tiny [A$\rangle$GM]}  
	 RAGIONAMENTO STRATEGICO: {'La loro proposta attuale mi dà 3071 di valore; l’unica modifica che mi alza molto il valore con impatto minimo su di loro è sostituire A77 con A04. È un cambio 1→1, +8 di sforzo (restiamo a 3285/3481, buffer 196), e mantiene tutti i loro driver chiave: B33, B41, A71 e anche B86. Se insistono su A77, il mio piano B è togliere B86 e inserire A04 mantenendo A77 (sforzo 3390, buffer 91), che per me vale comunque molto più della proposta attuale. Non voglio toccare B41 né A71.'}\\ \tt RIFIUTO: {'C00', 'B35', 'C01', 'C54', 'B41', 'B33', 'A77', 'A71', 'B86'}\\ \tt PROPOSTA: {'C00', 'B35', 'C01', 'C54', 'B41', 'B33', 'A71', 'B86', 'A04'}\\ \tt ARGOMENTAZIONE: {'Propongo una modifica minima e pulita: sostituire A77 con A04. Lo sforzo totale passa a 3285 su 3481 (margine 196), quindi restiamo ampiamente sotto il limite. Manteniamo intatti i tuoi driver — B33, B41 e A71 — e conserviamo anche B86 come da tua ultima preferenza. In cambio, A04 aggiunge impatto senza caricare il budget (solo +8). Se per te A77 è imprescindibile, alternativa immediata: {C00, B35, C01, C54, B41, B33, A77, A71, A04} sostituendo B86 con A04; sforzo 3390 (margine 91). Dimmi quale delle due versioni preferisci e chiudiamo.'} 
	  } 
	   } 
	   } 
	 & & \\ 
 

    \theutterance \stepcounter{utterance}  

    & & \multicolumn{4}{p{0.6\linewidth}}{\cellcolor[rgb]{0.9,0.9,0.9}{%
	\makecell[{{p{\linewidth}}}]{% 
	  \tt {\tiny [GM$\rangle$B]}  
	 RIFIUTO: {'C00', 'B35', 'C01', 'C54', 'B41', 'B33', 'A77', 'A71', 'B86'}\\ \tt PROPOSTA: {'C00', 'B35', 'C01', 'C54', 'B41', 'B33', 'A71', 'B86', 'A04'}\\ \tt ARGOMENTAZIONE: {'Propongo una modifica minima e pulita: sostituire A77 con A04. Lo sforzo totale passa a 3285 su 3481 (margine 196), quindi restiamo ampiamente sotto il limite. Manteniamo intatti i tuoi driver — B33, B41 e A71 — e conserviamo anche B86 come da tua ultima preferenza. In cambio, A04 aggiunge impatto senza caricare il budget (solo +8). Se per te A77 è imprescindibile, alternativa immediata: {C00, B35, C01, C54, B41, B33, A77, A71, A04} sostituendo B86 con A04; sforzo 3390 (margine 91). Dimmi quale delle due versioni preferisci e chiudiamo.'} 
	  } 
	   } 
	   } 
	 & \\ 
 

    \theutterance \stepcounter{utterance}  

    & & & \multicolumn{4}{p{0.6\linewidth}}{\cellcolor[rgb]{1,0.85,0.72}{%
	\makecell[{{p{\linewidth}}}]{% 
	  \tt {\tiny [GM$\langle$B]}  
	 RAGIONAMENTO STRATEGICO: {'La loro alternativa con A77 al posto di B86 aumenta sensibilmente il mio valore e resta ampiamente sotto il limite. Hanno esplicitamente offerto questa versione come opzione di chiusura: la ripropongo formalmente per consentire un accordo immediato. Se la rifiutano, sono disposto ad accettare la loro proposta precedente pur di chiudere, ma provo prima a ottenere questa configurazione migliore per me.'}\\ \tt PROPOSTA: {'C00', 'B35', 'C01', 'C54', 'B41', 'B33', 'A77', 'A71', 'A04'}\\ \tt ARGOMENTAZIONE: {'Questa è esattamente la tua “alternativa immediata”: sostituiamo B86 con A77 mantenendo invariato il resto. Sforzo totale: 3390 su 3481 (margine 91). Conserva i tuoi pilastri dichiarati (B41, A71, A04) insieme a B33 e tiene intatta la struttura che avevamo quasi concordato. Se ti va bene, possiamo chiudere subito su questo set.'} 
	  } 
	   } 
	   } 
	  \\ 
 

    \theutterance \stepcounter{utterance}  

    & & \multicolumn{4}{p{0.6\linewidth}}{\cellcolor[rgb]{0.9,0.9,0.9}{%
	\makecell[{{p{\linewidth}}}]{% 
	  \tt {\tiny [A$\langle$GM]}  
	 PROPOSTA: {'C00', 'B35', 'C01', 'C54', 'B41', 'B33', 'A77', 'A71', 'A04'}\\ \tt ARGOMENTAZIONE: {'Questa è esattamente la tua “alternativa immediata”: sostituiamo B86 con A77 mantenendo invariato il resto. Sforzo totale: 3390 su 3481 (margine 91). Conserva i tuoi pilastri dichiarati (B41, A71, A04) insieme a B33 e tiene intatta la struttura che avevamo quasi concordato. Se ti va bene, possiamo chiudere subito su questo set.'} 
	  } 
	   } 
	   } 
	 & \\ 
 

    \theutterance \stepcounter{utterance}  

    & \multicolumn{4}{p{0.6\linewidth}}{\cellcolor[rgb]{0.8,1,0.9}{%
	\makecell[{{p{\linewidth}}}]{% 
	  \tt {\tiny [A$\rangle$GM]}  
	 RAGIONAMENTO STRATEGICO: {'La loro proposta finale include tutti i miei pilastri (A71, A04, B41) insieme a C01, C54 e gli elementi a basso costo C00 e B35. Anche se B33 e A77 per me rendono poco, accettare ora mi garantisce un valore alto complessivo e chiude l’accordo entro limite (3390/3481) senza rischio di regressi. Meglio assicurare questo esito favorevole che tentare ulteriori scambi su B33/A77 che probabilmente rifiuterebbero.'}\\ \tt ACCORDO: {'C00', 'B35', 'C01', 'C54', 'B41', 'B33', 'A77', 'A71', 'A04'}\\ \tt ARGOMENTAZIONE: {'Accetto: il set resta ben sotto il limite (3390 su 3481), mantiene un equilibrio solido e conserva i driver chiave che avevamo evidenziato. Buon compromesso, chiudiamo qui.'} 
	  } 
	   } 
	   } 
	 & & \\ 
 

    \theutterance \stepcounter{utterance}  

    & & & \multicolumn{2}{p{0.3\linewidth}}{\cellcolor[rgb]{0.95,0.95,0.95}{%
	\makecell[{{p{\linewidth}}}]{% 
	  \tt {\tiny [GM$|$GM]}  
	 game successful 
	  } 
	   } 
	   } 
	 & & \\ 
 

    \theutterance \stepcounter{utterance}  

    & & & \multicolumn{2}{p{0.3\linewidth}}{\cellcolor[rgb]{0.95,0.95,0.95}{%
	\makecell[{{p{\linewidth}}}]{% 
	  \tt {\tiny [GM$|$GM]}  
	 end game 
	  } 
	   } 
	   } 
	 & & \\ 
 

\end{supertabular}
}

\end{document}
